%% Retained Hilbert-calculus proof after exact source restoration.
\begin{annexitem}{annex:p2.calc.2}{thm:p02-hilbert-deduction}{Deduction theorem for the Hilbert calculus}
\begin{proof}
First suppose
\(\Gamma\derives_{\mathsf H_{\mathsf{CPL}}}\varphi\to\psi\).
The same derivation is available from \(\Gamma\cup\{\varphi\}\); append
\(\varphi\) as a premise and then apply modus ponens.  Hence
\(\Gamma\cup\{\varphi\}\derives_{\mathsf H_{\mathsf{CPL}}}\psi\).

Conversely, let \(\theta_0,\ldots,\theta_n=\psi\) be a Hilbert derivation
from \(\Gamma\cup\{\varphi\}\).  Induct on \(i\) to derive
\(\varphi\to\theta_i\) from \(\Gamma\).

If \(\theta_i\in\Gamma\), or \(\theta_i\) is an axiom-schema instance,
derive \(\theta_i\), take the A1 instance
\(\theta_i\to(\varphi\to\theta_i)\), and apply modus ponens.  If
\(\theta_i=\varphi\), the following A1 and A2 instances derive
\(\varphi\to\varphi\):
\[
\begin{aligned}
 &\varphi\to((\varphi\to\varphi)\to\varphi),\\
 &\varphi\to(\varphi\to\varphi),\\
 &[\varphi\to((\varphi\to\varphi)\to\varphi)]
   \to([\varphi\to(\varphi\to\varphi)]
   \to(\varphi\to\varphi)).
\end{aligned}
\]
Apply modus ponens first to the first and third displayed lines, and then to
the result and the second line.

Finally suppose \(\theta_i\) follows by modus ponens from earlier
\(\eta\) and \(\eta\to\theta_i\).  The induction hypothesis gives
\(\varphi\to\eta\) and \(\varphi\to(\eta\to\theta_i)\).  Use the A2
instance
\[
 [\varphi\to(\eta\to\theta_i)]
 \to([\varphi\to\eta]\to[\varphi\to\theta_i])
\]
and modus ponens twice to obtain \(\varphi\to\theta_i\).  These exhaust
the ways a Hilbert line is licensed.  At \(i=n\) the result is
\(\Gamma\derives_{\mathsf H_{\mathsf{CPL}}}\varphi\to\psi\).
\end{proof}
\end{annexitem}

%% Proof of the propositional tableau adequacy theorem from the lemmas of
%% sec:p02-w02-p02-014-tableau-metatheory.
\begin{annexitem}{annex:p02-w02-p02-014-tableau-adequacy}{thm:p02-tableau-adequacy}{Tableau adequacy}
\begin{proof}
Tableaux, expandable occurrences and finished tableaux are as in
\cref{def:p02-w02-p02-014-input,def:p02-w02-p02-014-expandable}.  The first
sentence of the theorem is \cref{thm:p02-w02-p02-014-soundness}(3).  The
second is \cref{thm:p02-w02-p02-014-termination}(4): expanding unmarked
expandable occurrences, in any order, stops after at most
\(2^{\mu(B_\Gamma)+1}-2\) steps at a finished tableau, whose open branches are
saturated.  The valuation described in the third sentence is the branch
valuation \(v_B\); by \cref{thm:p02-w02-p02-014-truth} it satisfies every
formula on an open saturated branch \(B\), and by
\cref{thm:p02-w02-p02-014-invariants}(1) every member of \(\Gamma\) is on
\(B\).  The final sentence is \cref{thm:p02-w02-p02-014-completeness}(1),
read both ways: a finished tableau is closed exactly when \(\Gamma\) is
unsatisfiable, so it has an open branch exactly when \(\Gamma\) is
satisfiable.
\end{proof}
\end{annexitem}
